% OpenStax Calculus Volume 2 -- Section 2.9 Exercises: Calculus of the Hyperbolic Functions
% Exercises reproduced from OpenStax, licensed under CC BY 4.0.
% Source: https://openstax.org/books/calculus-volume-2/pages/2-9-calculus-of-the-hyperbolic-functions
\documentclass{exercises}
\title{OpenStax Calculus Volume 2}
\subtitle{Section 2.9 Exercises: Calculus of the Hyperbolic Functions}
\sourceurl{https://openstax.org/books/calculus-volume-2/pages/2-9-calculus-of-the-hyperbolic-functions}
\author{}
\date{}

\begin{document}
\maketitle

\begin{enumerate}[start=1]
\item \techrequired Find expressions for $\cosh x+\sinh x$ and $\cosh x-\sinh x$. Use a calculator to graph these functions and ensure your expression is correct.
\item From the definitions of $\cosh(x)$ and $\sinh(x)$, find their antiderivatives.
\item Show that $\cosh(x)$ and $\sinh(x)$ satisfy $y''=y$.
\item Use the quotient rule to verify that $\tanh(x)'=\operatorname{sech}^{2}(x)$.
\item Derive $\cosh^{2}(x)+\sinh^{2}(x)=\cosh(2x)$ from the definition.
\item Take the derivative of the previous expression to find an expression for $\sinh(2x)$.
\item Prove $\sinh(x+y)=\sinh(x)\cosh(y)+\cosh(x)\sinh(y)$ by changing the expression to exponentials.
\item Take the derivative of the previous expression to find an expression for $\cosh(x+y)$.
\end{enumerate}

For the following exercises, find the derivatives of the given functions and graph along with the function to ensure your answer is correct.

\begin{enumerate}[resume]
\item \techrequired $\cosh(3x+1)$
\item \techrequired $\sinh(x^{2})$
\item \techrequired $\frac{1}{\cosh(x)}$
\item \techrequired $\sinh(\ln(x))$
\item \techrequired $\cosh^{2}(x)+\sinh^{2}(x)$
\item \techrequired $\cosh^{2}(x)-\sinh^{2}(x)$
\item \techrequired $\tanh(\sqrt{x^{2}+1})$
\item \techrequired $\frac{1+\tanh(x)}{1-\tanh(x)}$
\item \techrequired $\sinh^{6}(x)$
\item \techrequired $\ln(\operatorname{sech}(x)+\tanh(x))$
\end{enumerate}

For the following exercises, find the antiderivatives for the given functions.

\begin{enumerate}[resume]
\item $\cosh(2x+1)$
\item $\tanh(3x+2)$
\item $x\cosh(x^{2})$
\item $3x^{3}\tanh(x^{4})$
\item $\cosh^{2}(x)\sinh(x)$
\item $\tanh^{2}(x)\operatorname{sech}^{2}(x)$
\item $\frac{\sinh(x)}{1+\cosh(x)}$
\item $\coth(x)$
\item $\cosh(x)+\sinh(x)$
\item $(\cosh(x)+\sinh(x))^{n}$
\end{enumerate}

For the following exercises, find the derivatives for the functions.

\begin{enumerate}[resume]
\item $\tanh^{-1}(4x)$
\item $\sinh^{-1}(x^{2})$
\item $\sinh^{-1}(\cosh(x))$
\item $\cosh^{-1}(x^{3})$
\item $\tanh^{-1}(\cos(x))$
\item $e^{\sinh^{-1}(x)}$
\item $\ln(\tanh^{-1}(x))$
\end{enumerate}

For the following exercises, find the antiderivatives for the functions.

\begin{enumerate}[resume]
\item $\int \frac{dx}{4-x^{2}}$
\item $\int \frac{dx}{a^{2}-x^{2}}$
\item $\int \frac{dx}{\sqrt{x^{2}+1}}$
\item $\int \frac{x\,dx}{\sqrt{x^{2}+1}}$
\item $\int-\frac{dx}{x\sqrt{1-x^{2}}}$
\item $\int \frac{e^{x}}{\sqrt{e^{2x}-1}}\,dx$
\item $\int-\frac{2x}{x^{4}-1}\,dx$
\end{enumerate}

For the following exercises, use the fact that a falling body with friction equal to velocity squared obeys the equation $dv/dt=g-v^{2}$.

\begin{enumerate}[resume]
\item Show that $v(t)=\sqrt{g}\tanh((\sqrt{g})t)$ satisfies this equation.
\item Derive the previous expression for $v(t)$ by integrating $\frac{dv}{g-v^{2}}=dt$.
\item \techrequired Estimate how far a body has fallen in $12$ seconds by finding the area underneath the curve of $v(t)$.
\end{enumerate}

For the following exercises, use this scenario: A cable hanging under its own weight has a slope $S=dy/dx$ that satisfies $dS/dx=c\sqrt{1+S^{2}}$. The constant $c$ is the ratio of cable density to tension.

\begin{enumerate}[resume]
\item Show that $S=\sinh(cx)$ satisfies this equation.
\item Integrate $dy/dx=\sinh(cx)$ to find the cable height $y(x)$ if $y(0)=1/c$.
\item Sketch the cable and determine how far down it sags at $x=0$.
\end{enumerate}

For the following exercises, solve each problem.

\begin{enumerate}[resume]
\item \techrequired A chain hangs from two posts $2$ m apart to form a catenary described by the equation $y=2\cosh(x/2)-1$. Find the slope of the catenary at the left fence post.
\item \techrequired A chain hangs from two posts four meters apart to form a catenary described by the equation $y=4\cosh(x/4)-3$. Find the total length of the catenary (arc length).
\item \techrequired A high-voltage power line is a catenary described by $y=10\cosh(x/10)$. Find the ratio of the area under the catenary to its arc length. What do you notice?
\item A telephone line is a catenary described by $y=a\cosh(x/a)$. Find the ratio of the area under the catenary to its arc length. Does this confirm your answer for the previous question?
\item Prove the formula for the derivative of $y=\sinh^{-1}(x)$ by differentiating $x=\sinh(y)$. (Hint: Use hyperbolic trigonometric identities.)
\item Prove the formula for the derivative of $y=\cosh^{-1}(x)$ by differentiating $x=\cosh(y)$. (Hint: Use hyperbolic trigonometric identities.)
\item Prove the formula for the derivative of $y=\operatorname{sech}^{-1}(x)$ by differentiating $x=\operatorname{sech}(y)$. (Hint: Use hyperbolic trigonometric identities.)
\item Prove that $(\cosh(x)+\sinh(x))^{n}=\cosh(nx)+\sinh(nx)$.
\item Prove the expression for $\sinh^{-1}(x)$. Multiply $x=\sinh(y)=(1/2)(e^{y}-e^{-y})$ by $2e^{y}$ and solve for $y$. Does your expression match the textbook?
\item Prove the expression for $\cosh^{-1}(x)$. Multiply $x=\cosh(y)=(1/2)(e^{y}+e^{-y})$ by $2e^{y}$ and solve for $y$. Does your expression match the textbook?
\end{enumerate}

\end{document}
